% HDC-RWKV walkthrough — every notation, every flow, on a hand-verifiable toy.
%
% Compile:  pdflatex architecture_walkthrough.tex
%
% Companion to architecture_hdc_rwkv.tex (the block diagram). This document
% takes a deliberately tiny instance — V=4, d=4, T=3, L=1 — and traces a
% training step and an inference step through every tensor in the model.
%
% Faithful to wozformer/models/hdc_rwkv.py and wozformer/ste.py.

\documentclass[11pt,a4paper]{article}

\usepackage[margin=22mm]{geometry}
\usepackage{amsmath,amssymb,amsthm}
\usepackage{mathtools}
\usepackage{booktabs}
\usepackage{array}
\usepackage{xcolor}
\usepackage{tcolorbox}
\usepackage{enumitem}
\usepackage{hyperref}
\usepackage{microtype}
\usepackage[T1]{fontenc}

\tcbuselibrary{breakable,skins}

\definecolor{cBipolar}{HTML}{2C5282}
\definecolor{cCont}{HTML}{B7791F}
\definecolor{cSampler}{HTML}{6B46C1}
\definecolor{cWarn}{HTML}{C53030}
\definecolor{cBoxBg}{HTML}{F7FAFC}

\newtcolorbox{notation}[1]{
  enhanced, breakable, colback=cBoxBg, colframe=cBipolar!70,
  fonttitle=\bfseries, title=#1, sharp corners=south,
}
\newtcolorbox{worked}[1]{
  enhanced, breakable, colback=cCont!4, colframe=cCont!80!black,
  fonttitle=\bfseries, title=#1, sharp corners=south,
}

% Matrices: row subscripts attach cleanly, so the "hv" tag is only added
% in the formal definition tables, not in every use.
\newcommand{\bV}{\mathbf{V}}
\newcommand{\bP}{\mathbf{P}}
\newcommand{\bVhv}{\mathbf{V}_{\!hv}}
\newcommand{\bPhv}{\mathbf{P}_{\!hv}}
\newcommand{\bm}{\mathbf{m}}
\newcommand{\bh}{\mathbf{h}}
\newcommand{\bs}{\mathbf{s}}
\newcommand{\bx}{\mathbf{x}}
\newcommand{\by}{\mathbf{y}}
\newcommand{\bz}{\mathbf{z}}
\newcommand{\bipset}{\{-1,+1\}}
\newcommand{\sg}{\,\text{sign}\,}
\newcommand{\stesign}{\,\widetilde{\text{sign}}\,}

\title{HDC-RWKV (nb12c): Every Notation, Every Flow\\
       \large A worked example for training and inference}
\author{Wozformer paper appendix}
\date{}

\begin{document}
\maketitle

\noindent This document accompanies the architecture diagram
(\texttt{architecture\_hdc\_rwkv.tex}). It defines every tensor and operator,
then traces a complete training step and a complete inference step through
the shipping model on a deliberately tiny instance ($V{=}4$, $d{=}4$, $T{=}3$,
$L{=}1$) so every number is verifiable by hand. The shipping configuration is
$V{=}256$, $d{=}256$, $L{=}1$, identical in every operation; only the dimensions
change.

\tableofcontents

\section{Conventions and symbols}

\begin{notation}{Indices and shapes}
\begin{tabular}{@{}lll@{}}
\toprule
Symbol & Meaning & Range / shape \\
\midrule
$B$       & batch size (number of independent sequences) & $\mathbb{N}_{>0}$ \\
$T$       & sequence length (block size) & $\mathbb{N}_{>0}$ \\
$V$       & vocabulary size (number of BPE tokens) & $\mathbb{N}_{>0}$ \\
$d$       & hypervector dimension & $\mathbb{N}_{>0}$ \\
$L$       & number of recurrent layers ($L{=}1$ in nb12c) & $\mathbb{N}_{>0}$ \\
$t$       & time step (token position) within a sequence & $0,1,\dots,T{-}1$ \\
$b$       & batch index & $0,1,\dots,B{-}1$ \\
$v$       & vocabulary index & $0,1,\dots,V{-}1$ \\
$i$       & hypervector channel index & $0,1,\dots,d{-}1$ \\
$x_t^{(b)}$ & input token id at position $t$ in sequence $b$ & $\{0,\dots,V{-}1\}$ \\
$y_t^{(b)}$ & target token (the answer for position $t$) & $\{0,\dots,V{-}1\}$ \\
\bottomrule
\end{tabular}
\end{notation}

\begin{notation}{Sets and spaces}
\begin{itemize}[leftmargin=18pt,itemsep=2pt]
  \item $\bipset^{d}$ --- the bipolar hypercube. Vectors of length $d$ whose
        entries are exactly $-1$ or $+1$. The natural state space of HDC.
  \item $(-1,+1)^{d}$ --- the open hypercube, image of $\tanh$. The
        \emph{continuous} state space we use for the recurrent activation.
  \item $\mathbb{R}^{V\times d}$ --- the continuous training-time storage for
        the bipolar tensors. At inference these tensors are reduced to a
        single bit per entry; at training they hold real-valued scratch values
        that get projected to $\bipset$ through STE on every forward pass.
\end{itemize}
\end{notation}

\section{The four learnable tensors}

The model has exactly four learnable objects. The first three are bipolar
(at inference); the fourth is a single scalar.

\begin{notation}{Learnable parameters}
\begin{tabular}{@{}lccp{60mm}@{}}
\toprule
Name (code) & Symbol & Shape (training) & Role \\
\midrule
\verb|vocab_hv_c|       & $\bVhv \in \mathbb{R}^{V\times d}$
   & $V \times d$ & Embedding table. Row $v$ is the \emph{token hypervector}
                    of vocabulary item $v$. Used through STE so the embedding
                    seen by the recurrence is $\sg(\bVhv) \in \bipset^{V\times d}$. \\
\verb|prototype_hv_c|   & $\bPhv \in \mathbb{R}^{V\times d}$
   & $V \times d$ & Output ``cleanup'' matrix. Row $v$ is the
                    \emph{prototype hypervector} for token $v$; the logit for
                    token $v$ is the inner product of the state with $\sg(\bPhv)_v$. \\
\verb|decay_masks_c[0]| & $\bm \in \mathbb{R}^{d}$
   & $d$ & Per-channel recurrence gate. After STE it lives in $\bipset^{d}$
          and tells each channel whether to keep $(+1)$ or flip $(-1)$ its
          previous state on each step. \\
\verb|log_temp|         & $\log\tau \in \mathbb{R}$
   & scalar & Output temperature, stored in log-space so $\tau > 0$ is
              automatic. Initialised to $\log\sqrt{d}$. \\
\bottomrule
\end{tabular}
\end{notation}

\begin{notation}{Deployment storage cost}
After training, only the \emph{sign} of each bipolar tensor matters. We pack
8 bits per byte:
\[
\text{bytes}_{\text{deploy}} \;=\; \underbrace{\tfrac{V\cdot d}{8}}_{\bV}
                           + \underbrace{\tfrac{V\cdot d}{8}}_{\bP}
                           + \underbrace{\tfrac{L\cdot d}{8}}_{\text{decay}}
                           + 2\;(\tau\text{ in fp16}).
\]
For nb12c ($V{=}256$, $d{=}256$, $L{=}1$): $8192 + 8192 + 32 + 2 \approx 16{,}418$ bytes;
the on-disk file adds a small header, totalling $16{,}434$ bytes.
\end{notation}

\section{Auxiliary operators}

Four operators appear in the forward pass; this section defines each.

\subsection{Element-wise sign with straight-through estimator}

\begin{notation}{$\stesign$ --- the operator that makes the model trainable}
The function $\stesign: \mathbb{R}^{n} \to \bipset^{n}$ has the
\emph{forward} value
\[
\stesign(x)_i \;=\; \begin{cases}
+1, & x_i \ge 0 \\
-1, & x_i < 0
\end{cases}
\]
and the \emph{backward} (subgradient) value
\[
\frac{\partial\,\stesign(x)_i}{\partial x_i} \;=\;
\begin{cases}
1, & |x_i| \le 1 \\
0, & |x_i| > 1
\end{cases}
\qquad\text{(``clipped STE'')}.
\]
In code this is realised by the algebraic identity
\[
\stesign(x) \;=\;
\underbrace{\sg(x).\text{detach}()}_{\text{forward value}} \;+\;
\underbrace{\text{clamp}_{[-1,1]}(x) - \text{clamp}_{[-1,1]}(x).\text{detach}()}_{\text{zero forward, identity gradient on }[-1,1]}.
\]
The detached terms have value but no gradient; the non-detached
$\text{clamp}_{[-1,1]}(x)$ supplies the gradient. See
\href{run:../../wozformer/ste.py}{\texttt{wozformer/ste.py}}.
\end{notation}

\subsection{Cyclic right-rotation $\rho^{t}$ (HDC binding)}

\begin{notation}{$\rho^{t}$ --- position binding via permutation}
For a hypervector $\bh \in \mathbb{R}^{d}$ and shift $t\in\mathbb{Z}$,
\[
\rho^{t}(\bh)_i \;=\; \bh_{\,(i - t)\bmod d}.
\]
Equivalent to \texttt{torch.roll(h, shifts=t, dims=-1)}. Key properties:
\begin{enumerate}[leftmargin=18pt,itemsep=2pt]
  \item \emph{Preserves bipolarity}: $\rho^{t}(\bipset^{d}) = \bipset^{d}$.
  \item \emph{Preserves quasi-orthogonality}: $\langle \rho^{a}(\bh),\rho^{b}(\bh)\rangle = d$ iff $a\equiv b\pmod d$,
        and is small otherwise.
  \item \emph{Composes additively}: $\rho^{a}\circ\rho^{b} = \rho^{a+b}$.
\end{enumerate}
Item (2) is why $\rho^{t}$ replaces a positional embedding table: positions
$t_1\ne t_2$ produce vectors that look uncorrelated to the recurrence, but the
operation is parameter-free.
\end{notation}

\subsection{Hyperbolic tangent}

Bounded continuous activation $\tanh: \mathbb{R}^{d} \to (-1,+1)^{d}$,
elementwise. Used to keep the recurrent state finite without binarising it
(finding F10 in \texttt{findings.md} shows binarising the state destroys
training).

\subsection{Cross-entropy loss}

For one position with logits $\bz \in \mathbb{R}^{V}$ and integer target
$y \in \{0,\dots,V{-}1\}$,
\[
\mathrm{CE}(\bz, y) \;=\; -\,\log\!\left(\frac{e^{\,z_y}}{\sum_{v=0}^{V-1} e^{\,z_v}}\right)
                       \;=\; \log\!\Bigl(\sum_v e^{z_v}\Bigr) - z_y.
\]
Averaged over all $B\cdot T$ positions.

\section{The forward pass --- definitions}

We give the forward pass first in symbolic form. \S\ref{sec:worked} then
executes every line on the toy instance.

\paragraph{Inputs.} A batch of token-id sequences
$\bx \in \{0,\dots,V{-}1\}^{B\times T}$ and (for training) targets
$\by \in \{0,\dots,V{-}1\}^{B\times T}$ where typically $y_t = x_{t+1}$.

\paragraph{Step 1 --- Materialise bipolar tensors.}
\begin{align}
\widetilde{\bV} \;&=\; \stesign(\bVhv) \;\in\; \bipset^{V\times d}, \\
\widetilde{\bP} \;&=\; \stesign(\bPhv) \;\in\; \bipset^{V\times d}, \\
\widetilde{\bm} \;&=\; \stesign(\bm) \;\in\; \bipset^{d}.
\end{align}

\paragraph{Step 2 --- Token lookup and position binding.}
For each position $t\in\{0,\dots,T-1\}$ and each batch element $b$,
\[
\bh_t^{(b)} \;=\; \rho^{\,t}\!\bigl(\,\widetilde{\bV}_{\,x_t^{(b)}}\,\bigr)
\quad\in\;\bipset^{d}.
\]
This produces a $(B,T,d)$ position-tagged bipolar tensor.

\paragraph{Step 3 --- Linear recurrence with continuous state.}
For each batch element, initialise $\bs_{-1}^{(b)} = \mathbf{0}\in\mathbb{R}^{d}$
and iterate
\[
\boxed{\;\;\bs_t^{(b)} \;=\; \tanh\!\Bigl(\,\widetilde{\bm}\odot \bs_{t-1}^{(b)} \;+\; \bh_t^{(b)}\,\Bigr)\;\;}
\qquad t=0,1,\dots,T-1,
\]
where $\odot$ is element-wise multiplication. Note: $\widetilde\bm \in \bipset^d$,
$\bs_{t-1} \in (-1,+1)^d$, $\bh_t \in \bipset^d$, so each channel of the
$\tanh$-input lies in $[-2,+2]$ and the state stays in $(-1,+1)^{d}$ by
construction.

\paragraph{Step 4 --- Prototype-similarity logits.}
With $\tau = e^{\log\tau}$,
\[
\bz_t^{(b)} \;=\; \frac{1}{\tau}\,\widetilde{\bP}\,\bs_t^{(b)} \;\in\;\mathbb{R}^{V}.
\]
Coordinate $v$ of $\bz_t^{(b)}$ is the inner product
$\langle \widetilde{\bP}_v, \bs_t^{(b)}\rangle / \tau$ --- the alignment of
the current state with prototype $v$.

\paragraph{Step 5 --- Loss (training only).}
\[
\mathcal{L} \;=\; \frac{1}{B\,T}\sum_{b=0}^{B-1}\sum_{t=0}^{T-1}
                  \mathrm{CE}\bigl(\bz_t^{(b)},\, y_t^{(b)}\bigr).
\]

\section{Worked training step on a toy instance}\label{sec:worked}

\paragraph{Toy configuration.} $V{=}4$, $d{=}4$, $T{=}3$, $L{=}1$, $B{=}1$.
Vocabulary: $\{\,\text{0:}\,\texttt{the},\;\text{1:}\,\texttt{cat},\;\text{2:}\,\texttt{sat},\;\text{3:}\,\texttt{mat}\,\}$.
Sequence: $\bx = (\texttt{the},\,\texttt{cat},\,\texttt{sat}) = (0,1,2)$, targets
$\by = (\texttt{cat},\,\texttt{sat},\,\texttt{mat}) = (1,2,3)$ (next-token).

\begin{worked}{Initial parameter values}
We hand-pick small integers so every step is checkable by mental arithmetic.
\[
\bV =
\begin{pmatrix}
+0.4 & -0.2 & +0.7 & -0.1 \\
-0.3 & +0.6 & -0.5 & +0.2 \\
+0.1 & +0.8 & -0.4 & -0.6 \\
+0.5 & -0.7 & +0.3 & +0.9
\end{pmatrix}
\;\;\Rightarrow\;\;
\widetilde\bV = \stesign(\bV) =
\begin{pmatrix}
+1 & -1 & +1 & -1 \\
-1 & +1 & -1 & +1 \\
+1 & +1 & -1 & -1 \\
+1 & -1 & +1 & +1
\end{pmatrix}.
\]
\[
\bP \;\Rightarrow\;
\widetilde\bP =
\begin{pmatrix}
+1 & +1 & -1 & -1 \\
-1 & +1 & +1 & -1 \\
+1 & -1 & +1 & -1 \\
-1 & -1 & -1 & +1
\end{pmatrix},\qquad
\bm = \bigl(+0.5,\;+0.5,\;+0.5,\;+0.5\bigr) \;\Rightarrow\;
\widetilde\bm = (+1,+1,+1,+1).
\]
\[
\log\tau = \tfrac{1}{2}\log d = \tfrac{1}{2}\log 4 \approx 0.693,
\qquad \tau = \sqrt{d} = 2.
\]
\end{worked}

\subsection{Step 1 --- Token lookup and rotation}

\begin{worked}{Position binding by $\rho^{t}$}
\textbf{$t=0$, token \texttt{the} (id 0).}\quad
$\widetilde\bV_0 = (+1,-1,+1,-1)$.\quad
$\rho^{0}$ is the identity:
$\bh_0 = (+1,-1,+1,-1).$

\medskip
\textbf{$t=1$, token \texttt{cat} (id 1).}\quad
$\widetilde\bV_1 = (-1,+1,-1,+1)$. Right-rotate by 1:
$(-1,+1,-1,+1)\!\overset{\rho^{1}}{\longrightarrow}\!(+1,-1,+1,-1)$.
So $\bh_1 = (+1,-1,+1,-1).$\footnote{For $d{=}4$, rotation by $1$ takes
$(a_0,a_1,a_2,a_3)\mapsto(a_3,a_0,a_1,a_2)$.}

\medskip
\textbf{$t=2$, token \texttt{sat} (id 2).}\quad
$\widetilde\bV_2 = (+1,+1,-1,-1)$. Right-rotate by 2:
$(+1,+1,-1,-1)\!\overset{\rho^{2}}{\longrightarrow}\!(-1,-1,+1,+1)$.
So $\bh_2 = (-1,-1,+1,+1).$
\end{worked}

\subsection{Step 2 --- Recurrence}

With $\widetilde\bm = (+1,+1,+1,+1)$ the recurrence simplifies to
$\bs_t = \tanh(\bs_{t-1} + \bh_t)$.

\begin{worked}{State evolution}
$\bs_{-1} = (0,0,0,0).$

\smallskip
\textbf{$t=0$:} update $= \bh_0 = (+1,-1,+1,-1)$.\\
$\bs_0 = \tanh(+1,-1,+1,-1) = (+0.7616,\,-0.7616,\,+0.7616,\,-0.7616).$

\smallskip
\textbf{$t=1$:} $\bh_1 = (+1,-1,+1,-1)$, so
$\bs_0 + \bh_1 = (0.7616{+}1,\;-0.7616{-}1,\;0.7616{+}1,\;-0.7616{-}1)
= (+1.7616,\,-1.7616,\,+1.7616,\,-1.7616).$\\
$\bs_1 = \tanh(\cdot) = (+0.9434,\,-0.9434,\,+0.9434,\,-0.9434).$

\smallskip
\textbf{$t=2$:} update $= \bs_1 + \bh_2 = (0.9434{-}1,\,-0.9434{-}1,\,0.9434{+}1,\,-0.9434{+}1)
= (-0.0566,\,-1.9434,\,+1.9434,\,+0.0566).$\\
$\bs_2 = \tanh(\cdot) = (-0.0566,\,-0.9601,\,+0.9601,\,+0.0566).$\footnote{$\tanh(0.0566)\approx 0.0565$; $\tanh(1.9434)\approx 0.9601$.}
\end{worked}

\noindent Already we can see what the recurrence does: channels for which
$\bh_t$ has the same sign across positions accumulate strongly; channels
where $\bh_t$ flips sign cancel. With $\widetilde\bm=+1$ everywhere this is
``always remember''; in a trained model some $m_i$ are $-1$, giving
``flip every step''.

\subsection{Step 3 --- Logits and loss at the final position}

At $t=2$ the target token is $y_2 = 3$ (\texttt{mat}).
We compute $\bz_2 = \widetilde\bP\,\bs_2 / \tau$ with $\tau = 2$.

\begin{worked}{Per-token logits at t=2}
$\bs_2 = (-0.0566,\,-0.9601,\,+0.9601,\,+0.0566).$
\begin{align*}
\langle\widetilde\bP_0,\bs_2\rangle &= (+1)(-0.0566)+(+1)(-0.9601)+(-1)(+0.9601)+(-1)(+0.0566) \\
 &= -0.0566 - 0.9601 - 0.9601 - 0.0566 \;=\; -2.0334, \\
\langle\widetilde\bP_1,\bs_2\rangle &= -(-0.0566)+(-0.9601)+(+0.9601)-(+0.0566) \;=\; 0.0000, \\
\langle\widetilde\bP_2,\bs_2\rangle &= (-0.0566)-(-0.9601)+(+0.9601)-(+0.0566) \;=\; +1.8070, \\
\langle\widetilde\bP_3,\bs_2\rangle &= -(-0.0566)-(-0.9601)-(+0.9601)+(+0.0566) \;=\; +0.1132.
\end{align*}
Dividing by $\tau = 2$:
\[
\bz_2 \;=\; (-1.0167,\;0.0000,\;+0.9035,\;+0.0566).
\]
\end{worked}

\begin{worked}{Cross-entropy loss at t=2}
Target is $y_2 = 3$. Apply softmax to $\bz_2$:
\[
e^{\bz_2} = (0.3617,\,1.0000,\,2.4683,\,1.0583),\qquad \sum = 4.8883.
\]
\[
p_3 = \frac{e^{z_{2,3}}}{\sum} = \frac{1.0583}{4.8883} \approx 0.2165,
\qquad \mathrm{CE} = -\log(0.2165) \approx 1.530.
\]
A correctly trained model would have made $p_3$ much larger than the other
three; here all four logits are within $\approx 2$ of one another, so the
loss is large.
\end{worked}

\subsection{Step 4 --- Backward pass through STE}

The graph of dependencies for $\mathcal{L}$ on $\bV_0$ (the row of
\verb|vocab_hv_c| corresponding to token $0$, \texttt{the}) is:
\[
\bV_0 \;\xrightarrow{\stesign}\; \widetilde\bV_0
      \;\xrightarrow{\rho^{0}}\; \bh_0
      \;\xrightarrow{\text{recurrence}}\; \bs_0 \to \bs_1 \to \bs_2
      \;\xrightarrow{\widetilde\bP\,\cdot\,/\tau}\; \bz_2
      \;\xrightarrow{\mathrm{CE}}\; \mathcal{L}.
\]
The only step that is not analytically differentiable is $\stesign$.
By definition (clipped STE),
\[
\Bigl[\nabla_{\bV_0} \mathcal{L}\Bigr]_{i} \;=\;
\Bigl[\nabla_{\widetilde\bV_0} \mathcal{L}\Bigr]_{i} \cdot \mathbf{1}\bigl[|\bV_{0,i}| \le 1\bigr].
\]
Concretely, $\bV_{0} = (+0.4,-0.2,+0.7,-0.1)$ all have $|\cdot|\le 1$, so the
indicator is $1$ for every channel and the gradient is passed through
unchanged. If at a later training step some $\bV_{0,i}$ drifts past $\pm 1$,
the gradient for that channel is gated to $0$ --- preventing runaway away
from the bipolar manifold.

\begin{worked}{What the optimiser actually updates}
Adam with $\eta = 3\!\cdot\!10^{-3}$ takes
\[
\bV \;\leftarrow\; \bV \;-\; \eta\,\widehat{\nabla_{\bV}\mathcal{L}},
\]
where the hat is Adam's bias-corrected step. The continuous storage $\bV$
drifts slowly; its \emph{sign}, $\widetilde\bV$, is what the recurrence sees
on the next forward and is what gets exported as the 1-bit-per-cell deployment
artifact. Sign flips of $\bV_{v,i}$ are the only events that change behaviour.
\end{worked}

\subsection{Recap of the full training step}

\begin{enumerate}[leftmargin=18pt,itemsep=2pt]
  \item Forward materialises $\widetilde\bV, \widetilde\bP, \widetilde\bm$ via $\stesign$.
  \item Tokens are looked up and rotated, producing $\bh_0,\bh_1,\bh_2$.
  \item The recurrence produces continuous states $\bs_0,\bs_1,\bs_2$.
  \item Each $\bs_t$ is dotted with $\widetilde\bP$ and divided by $\tau$ to get logits.
  \item Cross-entropy against $\by$ is the loss.
  \item Backward: gradients flow through $\tanh$, the inner product, and the
        clipped STE; only continuous shadows ($\bV,\bP,\bm,\log\tau$) get
        gradient updates.
\end{enumerate}

\section{Worked inference (autoregressive sampling)}

Inference uses the same forward pass, with three differences:
\begin{enumerate}[leftmargin=18pt,itemsep=2pt]
  \item No targets, so no loss; we stop at logits.
  \item The bipolar tensors come from \verb|.sign()|, not \verb|ste_sign|;
        the gradient path does not exist. The values are identical.
  \item Each new token is sampled and fed back as $x_{t+1}$; the loop continues
        for \texttt{max\_new\_tokens} steps.
\end{enumerate}

\paragraph{Toy setting.} Same parameters as \S\ref{sec:worked}. Prompt: a single
token, \texttt{the} (id 0). Generate one new token using
\emph{argmax} (deterministic, equivalent to top-1).

\begin{worked}{Forward on the single-token prompt}
$\bs_{-1} = (0,0,0,0)$, $\bh_0 = (+1,-1,+1,-1)$.\\
$\bs_0 = \tanh(\bh_0) = (+0.7616,\,-0.7616,\,+0.7616,\,-0.7616).$
\end{worked}

\begin{worked}{Logits at the first generation step}
\begin{align*}
\langle\widetilde\bP_0,\bs_0\rangle &= (+1)(+0.7616) + (+1)(-0.7616) + (-1)(+0.7616) + (-1)(-0.7616) = 0.0000, \\
\langle\widetilde\bP_1,\bs_0\rangle &= -0.7616 - 0.7616 + 0.7616 + 0.7616 = 0.0000, \\
\langle\widetilde\bP_2,\bs_0\rangle &= +0.7616 + 0.7616 + 0.7616 + 0.7616 = +3.0464, \\
\langle\widetilde\bP_3,\bs_0\rangle &= -0.7616 + 0.7616 - 0.7616 - 0.7616 = -1.5232.
\end{align*}
$\bz_0 = (0,\,0,\,+1.5232,\,-0.7616)$ after $\tau = 2$.
\end{worked}

\begin{worked}{Sampling decision}
Argmax over $\bz_0$ picks index $2 \Rightarrow$ token \texttt{sat}.

\medskip
Top-$k$ with $k{=}2$, temperature $T_{\!s}{=}0.7$ would instead form
$p \propto \exp(\bz_0 / T_{\!s})$ restricted to the top-2 indices
$\{0,2\}$ (the others' probabilities are set to $0$ before normalisation):
\[
\exp(0/0.7) = 1,\quad \exp(1.5232/0.7) \approx 8.846,
\qquad p_0 \approx 0.102,\;\; p_2 \approx 0.898.
\]
With a random draw $u \sim \mathcal{U}[0,1]$ from a seeded RNG, we emit
\texttt{sat} unless $u < 0.102$, in which case we emit \texttt{the}.
\end{worked}

\paragraph{Continuing the loop.} The sampled token is appended to the input
and the forward repeats with $T'\!=\!2$ (positions $0$ and $1$). Because the
recurrence is left-to-right and stateful, an efficient implementation caches
$\bs_{T'-1}$ and runs the recurrence one step. The 6502 firmware does exactly
this: it keeps a 256-byte state register and updates it once per emitted
token.

\section{What the 6502 actually computes}

At deployment the math reduces to integer bit operations.

\begin{notation}{Bipolar dot product as XOR + popcount}
For $a, b \in \bipset^{d}$, encode $-1\mapsto 0$ and $+1\mapsto 1$, giving
$\bar a,\bar b\in\{0,1\}^{d}$. Then
\[
\langle a, b\rangle \;=\; d - 2\cdot\mathrm{popcount}(\bar a \oplus \bar b),
\]
where $\oplus$ is bitwise XOR. The mapping is exact: agreements add $+1$
each, disagreements add $-1$ each.
\end{notation}

\begin{notation}{Inference loop in bit-arithmetic terms}
At each emitted token:
\begin{enumerate}[leftmargin=18pt,itemsep=2pt]
  \item \textbf{Token lookup:} fetch the $d/8$ bytes of $\widetilde\bV_{x_t}$ from EEPROM.
  \item \textbf{Rotation:} cyclic shift those bytes by $t$ bits.
  \item \textbf{Recurrence:} for each of $d$ channels, compute
        $\widetilde m_i\cdot s_{i,t-1} + h_{i,t}$ in fixed-point and apply a
        $\tanh$ LUT.
  \item \textbf{Logits:} for each of $V$ rows of $\widetilde\bP$, XOR with the
        thresholded state and popcount; convert to a fixed-point logit.
  \item \textbf{Sample:} softmax + top-$k$ + seeded PRNG.
\end{enumerate}
The two storage-heavy tables ($\widetilde\bV, \widetilde\bP$) account for
$16{,}384$ of the $16{,}434$ bytes; everything else is scratch.
\end{notation}

\section{Cross-reference to the source}

Every equation above corresponds to a specific line range:

\begin{tabular}{@{}lll@{}}
\toprule
Concept & File & Lines \\
\midrule
$\stesign$ definition           & \texttt{wozformer/ste.py}                   & 15--24 \\
Parameter declarations          & \texttt{wozformer/models/hdc\_rwkv.py}      & 30--36 \\
Forward (training, STE)         & \texttt{wozformer/models/hdc\_rwkv.py}      & 39--74 \\
Forward (deployment, \verb|.sign()|) & \texttt{wozformer/models/hdc\_rwkv.py} & 78--111 \\
Deployment byte counter         & \texttt{wozformer/models/hdc\_rwkv.py}      & 114--117 \\
\bottomrule
\end{tabular}

\bigskip
\noindent
The block diagram in \texttt{architecture\_hdc\_rwkv.tex} visualises the
same flow at the shipping dimensions $V{=}256$, $d{=}256$. Only the magnitudes
change; the math is identical.

\end{document}
